\documentclass[10pt,twocolumn]{article}

% ── Packages ──────────────────────────────────────────────────
\usepackage[margin=1in]{geometry}
\usepackage{amsmath,amssymb}
\usepackage{graphicx}
\usepackage{booktabs}
\usepackage{hyperref}
\usepackage{cite}
\usepackage{caption}
\usepackage{subcaption}
\usepackage{xcolor}
\usepackage{microtype}

% ── Title ─────────────────────────────────────────────────────
\title{\textbf{Latent Representations of Diffusion Models\\for Deepfake Detection}\\
  \large A Controlled Ablation Study}

\author{
  Loïc Huang\thanks{Loic.Huang@eurecom.fr} \and
  Wenyi Zeng\thanks{Wenyi.Zeng@eurecom.fr} \\[0.3em]
  \normalsize EURECOM, Sophia Antipolis, France \\
  \normalsize Supervisors: Jean-Luc Dugelay \& Alexandre Libourel \\
  \normalsize Semester Research Project, Spring 2026
}

\date{}

\begin{document}

\maketitle

% ── Abstract ──────────────────────────────────────────────────
\begin{abstract}
% TODO: 150-200 words summarising the problem, approach, key results.
% H1 validated (0.87 vs 0.78), generalisation wall, fine-tuning trap,
% denoising residual +8 pts cross-dataset.
\end{abstract}

% ── 1. Introduction ───────────────────────────────────────────
\section{Introduction}
% TODO:
% - Context: deepfakes and societal stakes
% - Detection challenge: generalisation across datasets
% - Gap: existing methods (Wang, Ricker) rely on full diffusion pipeline
% - Our question: is the VAE encoder alone sufficient? (H1)
% - Contributions: controlled ablation, fine-tuning study, denoising residual
% - Paper structure

% ── 2. Related Work ───────────────────────────────────────────
\section{Related Work}

\subsection{Deepfake Detection}
% TODO: CNN-based detectors, ViT-based, limitations on generalisation

\subsection{Diffusion Models as Detectors}
% TODO: Wang & Kalogeiton (ECCV 2024), Ricker et al. AEROBLADE (CVPR 2024)

\subsection{Latent Representations}
% TODO: VAE encoder properties, LAION pretraining, generative priors

% ── 3. Method ─────────────────────────────────────────────────
\section{Method}

\subsection{Hypothesis}
% TODO: H1 — frozen VAE encoder sufficient for detection

\subsection{Controlled Comparison Pipeline}
% TODO: LDM VAE encoder vs ResNet-50, shared MLP, same protocol

\subsection{Variant A: Fine-tuning the Encoder}
% TODO: Trainable VAE encoder vs trainable ResNet-50 backbone, same MLP head

\subsection{Variant B: The Denoising Residual}
% TODO: |z0 - z_hat0|, Try 1 (1-step), Try 2 (20-step), inspiration from Wang

% ── 4. Experimental Setup ─────────────────────────────────────
\section{Experimental Setup}

\subsection{Datasets}
% TODO: FaceForensics++ (DeepFakes, c23), train/val/test split
%       Celeb-DF-v2 (cross-dataset benchmark, strictly unseen)

\subsection{Evaluation Protocol}
% TODO: AUROC, video-level, 5 frames per video
%       95% bootstrap CI, 2000 iterations, cluster resampling

\subsection{Implementation Details}
% TODO: SD 1.4 VAE, ResNet-50 ImageNet1K-V2, MLP architecture,
%       learning rates, batch size, epochs, hardware

% ── 5. Results ────────────────────────────────────────────────
\section{Results}

\subsection{Frozen Encoders (H1)}
% TODO: Table — LDM 0.87 vs ResNet 0.78 FF++, both ~0.52 Celeb-DF
%       Confidence intervals, H1 validated in-domain

\subsection{Fine-tuning Results}
% TODO: Table — LDM FT 0.98/0.50, ResNet FT 0.97/0.70
%       The Specialisation Trap

\subsection{Denoising Residual}
% TODO: Table — Try1 0.56/0.50, Try2 0.79/0.60, Wang ref 0.70
%       Effect of reconstruction depth

% ── 6. Analysis ───────────────────────────────────────────────
\section{Analysis and Discussion}

\subsection{Why Does the LDM Encoder Generalise Better In-Domain?}
% TODO: Generative prior, LAION pretraining, richer feature space

\subsection{The Specialisation Trap}
% TODO: LDM FT overfits FF++ latent patterns, loses generative prior
%       ResNet discriminative bias more transferable

\subsection{Limits of the Denoising Residual}
% TODO: Gap vs Wang: |z0-z_hat0| vs multi-step eps_theta
%       Compute cost of Try 2

% ── 7. Conclusion ─────────────────────────────────────────────
\section{Conclusion}
% TODO: 4 key takeaways, future work (LoRA, training on all FF++ manipulations)

% ── References ────────────────────────────────────────────────
\bibliographystyle{plain}
\bibliography{references}

% Manually if no .bib file yet:
% \begin{thebibliography}{9}
% \bibitem{wang2024} Wang \& Kalogeiton, ECCV 2024.
% \bibitem{ricker2024} Ricker et al., CVPR 2024.
% \bibitem{rossler2019} Rössler et al., FaceForensics++, ICCV 2019.
% \bibitem{li2020} Li et al., Celeb-DF, CVPR 2020.
% \end{thebibliography}

\end{document}
